% ============================================================
%  method.tex — 方法:可复现细节充足,读者能照做
%  纪律:公式符号一致;架构/损失/超参齐全;区分「方法本身」与「实现选择」。
% ============================================================

\section{Method}
\label{sec:method}

% [问题形式化] 记号定义(POMDP / 策略 π(a|o) 等),符号全篇一致。
\subsection{Problem Formulation}
We consider \ldots. Let $o_t$ denote \ldots and $a_t$ denote \ldots.
The goal is to learn a policy $\pi_\theta(a_t \mid o_t, \ell)$ conditioned on
language instruction $\ell$ \ldots

% [方法总览] 配架构图(见 ⑦ figure-builder, diagram 类)。
\subsection{Overview}
Figure~\ref{fig:pipeline} illustrates MethodName, which consists of \ldots
\begin{figure}[t]
  \centering
  \includegraphics[width=\linewidth]{pipeline}
  \caption{Overview of MethodName. \ldots (caption 须自解释)}
  \label{fig:pipeline}
\end{figure}

% [核心组件] 逐个讲清提出的模块;每个组件对应后文一个消融。
\subsection{Component 1: <名称>}
\ldots
\begin{equation}
  \mathcal{L} = \ldots
  \label{eq:loss}
\end{equation}

\subsection{Component 2: <名称>}
\ldots

% [训练/推理] 算法伪代码(algorithm 环境);关键超参见附录,不散落正文。
\subsection{Training and Inference}
\begin{algorithm}[t]
\caption{MethodName Training}
\begin{algorithmic}[1]
\STATE \ldots
\end{algorithmic}
\end{algorithm}

% [实现选择 vs 方法本质] 明确哪些是可替换的工程选择,避免过度声称普适性。
\paragraph{Design choices.} We emphasize that \ldots is an implementation choice,
whereas \ldots is central to the method.
